Open vision-language models (VLMs) for Thai OCR now read most of a page correctly, but Thai vowels and tone marks are small, and one mark can change a word. This paper asks where the marks are lost, in addition to how many. We introduce a mark score that ignores reading order, penalizes surplus text, and gives no credit to the copies produced by a repetition loop. On the ThaiOCRBench calibration split, Typhoon OCR 1.5 (2B) writes a wrong tone mark in only 0.65\% of the places where it reads the base consonant correctly. Given the correct preceding text, it chooses the right tone mark 98.8\% of the time even without the image. Most of the marks it loses are never read: lines are absent, spans are missing, or a loop consumes the output. Among training-free decoders, stopping at a detected loop leaves every non-looping output unchanged and raises mark F1 by 0.8 points for Typhoon and by 10.8 points for a 0.9B Thai model. A repetition penalty lowers F1 on pages that do not loop, and subtracting the language prior removes correct tone marks. \todo{Results for other models and a locked-split confirmation are pending.} All results are preliminary. They come from one benchmark and a calibration split, and some rules were written after the data had been seen.
